\section{Следующие инженерные гипотезы}

Из текущих отрицательных результатов следует важная развилка: не всякая работа с
ID маршрутов одинаково полезна.

\subsection{Простая глобальная сортировка ID}

Самая дешёвая идея --- просто перенумеровать маршруты по частоте или по близости.
Это может слегка улучшить locality lookup'ов, но вряд ли решит основную проблему,
поскольку текущий codebook и так мал и уже не является главным bottleneck.
Иными словами, простая сортировка ID выглядит как полезный housekeeping-шаг, но
не как основной источник speed win.

\subsection{Tile-local grouping и локальные palette}

Более сильная гипотеза --- отказаться от чисто глобального словаря во время
исполнения и перейти к tile-local представлению. Тогда для каждого weight-tile
можно хранить небольшой локальный palette маршрутов и компактные локальные
индексы. Kernel один раз загружает этот palette в shared memory и затем
многократно переиспользует его по всей tile.

Такой подход атакует уже не codebook locality, а главный источник перерасхода ---
bandwidth матрицы ID и per-element decode path.

\subsection{Runtime-aware route distillation как аналог bitdistill}

Наиболее интересная исследовательская гипотеза --- ввести для DRL собственный
аналог bitdistill, но ориентированный не только на качество, а на runtime-
структуру.

Идея состоит в том, чтобы обучать route-student под руководством dense teacher
или уже качественного routed-teacher, но добавлять к quality-loss также
структурные регуляризаторы:

\begin{itemize}[leftmargin=1.5em]
\item уменьшение числа уникальных маршрутов внутри tile;
\item уменьшение route-entropy внутри tile;
\item поощрение стабильного повторного использования маршрутов в соседних
позициях;
\item ограничение размера локального palette, удобного для shared-memory kernel.
\end{itemize}

Такой distillation меняет не только численные веса, но и саму геометрию route-
языка так, чтобы он был аппаратно удобнее. Это выглядит существенно более
сильным направлением, чем простая постфактум сортировка ID.